\documentclass[../DoAn.tex]{subfiles}

\begin{document}

\begin{center}
    {\bfseries\fontsize{18pt}{22pt}\selectfont ABSTRACT}
\end{center}
\vspace{1.5cm}

Dormitory management at many Vietnamese universities still relies largely on
manual procedures that are labour-intensive, error-prone, and lack transparency
both in allocating rooms and in explaining the outcomes to students. Motivated by
this situation, the thesis develops \textbf{eDorm} --- an integrated web and
mobile platform for end-to-end online dormitory management, registration, and
room allocation, aiming to automate the process, make it transparent, and improve
the digital experience for both administrators and students. Methodologically, the
system performs automatic room allocation based on multi-criteria priority scores,
provides a sandbox simulation environment that allows administrators to experiment
and roll back changes with a single click, and forecasts accommodation demand
under four scenarios; the user interfaces are built with \texttt{React Native}
together with \texttt{Expo} and a Progressive Web App, while the residential space
is visualised through a \texttt{CesiumJS} three-dimensional map and a
\texttt{Three.js} virtual reality room tour. Experimental results show that the
system operates stably at around 300 concurrent users and was tested on a dataset
of more than 1{,}700 students across seven dormitory buildings. With these results,
\textbf{eDorm} demonstrates its feasibility and can be deployed in practice at
Vietnamese universities to comprehensively modernise dormitory management.

\vspace{2cm}
\begin{flushright}
    Student\\[4pt]
    {\itshape PHAN ĐỨC QUYỀN}
\end{flushright}

\end{document}
